%%% FIELD proof-flat-face-exact-value
Put \(a=1-\alpha\).  Since \(\alpha\in(0,1/4)\), we have
\(a\in(3/4,1)\).  Fix an index \(n\) such that
\(\kappa_n\le 1/16\).  By
\(\text{thm:integrated-and-bump-lower-frontier}\),
\[
 L_n^\star(\kappa_n)\ge G_\alpha(\tau_n).
\]
By \(\text{def:integrated-slope-lower-function}\),
\[
 G_\alpha(\tau_n)
 =2\int_0^1 [a-4\sqrt{\tau_n t}]_+\,dt.
\]
For every \(t\in[0,1]\), one has
\(\sqrt{\tau_n t}\le \sqrt{\tau_n}\).  Because the map
\(x\mapsto[x]_+\) is nondecreasing,
\[
 [a-4\sqrt{\tau_n t}]_+
 \ge [a-4\sqrt{\tau_n}]_+.
\]
Integrating this pointwise inequality over \([0,1]\) gives
\[
 L_n^\star(\kappa_n)
 \ge G_\alpha(\tau_n)
 \ge 2[a-4\sqrt{\tau_n}]_+
 =2[1-\alpha-4\sqrt{\tau_n}]_+.
\]

For the upper bound, define the measurable sample rule
\[
 I_n^0(O_1,\ldots,O_n)
 =
 \begin{cases}
   \Theta,& X_1\le 1-\alpha,\\
   \{0\},& X_1>1-\alpha.
 \end{cases}
\]
Its two possible values are connected subsets of
\(\Theta=[-1,1]\), with lengths \(2\) and \(0\), respectively.
For every \(P\in\mathcal P_n\), the parameter-space restriction gives
\(\theta_P\in\Theta\).  Therefore
\[
 \{X_1\le1-\alpha\}
 \subseteq
 \{\theta_P\in I_n^0(O_1,\ldots,O_n)\}.
\]
By \(\text{ass:uniform-running-variable}\),
\[
 P\{\theta_P\in I_n^0(O_1,\ldots,O_n)\}
 \ge P\{X_1\le1-\alpha\}
 =1-\alpha.
\]
Thus \(I_n^0\) obeys the uniform honesty constraint in
\(\text{def:minimax-length}\).  The same uniform-design assumption
also yields, for every \(P\in\mathcal P_n\),
\[
 E_P\bigl[|I_n^0(O_1,\ldots,O_n)|\bigr]
 =2P\{X_1\le1-\alpha\}
 =2(1-\alpha).
\]
Using this single data-driven rule as a competitor in the infimum in
\(\text{def:minimax-length}\) proves
\[
 L_n^\star(\kappa_n)\le2(1-\alpha).
\]
This establishes both finite-sample inequalities.

Now suppose \(\tau_n\to0\).  Since \(n\ge1\) and
\(\tau_n=n\kappa_n^3\),
\[
 0<\kappa_n^3\le\tau_n,
\]
so \(\kappa_n\to0\).  Consequently \(\kappa_n\le1/16\)
eventually, and the preceding bounds apply for all sufficiently large
\(n\).  Their lower endpoint converges to \(2(1-\alpha)\), while their
upper endpoint equals \(2(1-\alpha)\).  The squeeze theorem therefore
gives
\[
 L_n^\star(\kappa_n)\longrightarrow2(1-\alpha).
\]

Finally, fix \(s>1\) and suppose that \(\lambda_n\) remains in a
compact subset of \((0,\infty)\).  In particular, \(\lambda_n\) is
bounded above.  We first show that \(\kappa_n\to0\).  Otherwise there
would be \(\epsilon>0\) and a subsequence, indexed again by \(n\), on
which \(\kappa_n\ge\epsilon\).  The declared strength range
\(\kappa_n\le1/8\) then permits \(0<\epsilon\le1/8\).  Since
\(x\mapsto\log(e/x)\) is decreasing, the definition of \(\lambda_n\)
implies along that subsequence
\[
 \lambda_n
 =\frac{n\kappa_n^{2+1/s}}{\log(e/\kappa_n)}
 \ge
 \frac{n\epsilon^{2+1/s}}{\log(e/\epsilon)}
 \longrightarrow\infty,
\]
contradicting boundedness.  Hence \(\kappa_n\to0\), and in particular
\(\kappa_n\le1/16\) eventually.  From the defining identity
\[
 \lambda_n
 =\frac{\tau_n\kappa_n^{1/s-1}}{\log(e/\kappa_n)}
\]
we obtain
\[
 \tau_n
 =\lambda_n\kappa_n^{1-1/s}\log(e/\kappa_n).
\]
Because \(s>1\), the exponent \(1-1/s\) is positive.  The elementary
limit
\[
 x^{1-1/s}\log(e/x)\longrightarrow0
 \qquad\text{as }x\downarrow0
\]
and the upper boundedness of \(\lambda_n\) therefore imply
\(\tau_n\to0\).  Applying the preceding squeeze on the eventual range
\(\kappa_n\le1/16\) proves
\[
 L_n^\star(\kappa_n)\longrightarrow2(1-\alpha),
\]
as required.
